\documentclass[nonacm]{acmart}

\AtBeginDocument{%
  \providecommand\BibTeX{{%
    Bib\TeX}}}

\setcopyright{none}

\usepackage{array}
\usepackage{mathtools}
\usepackage{stmaryrd}
\usepackage{mathpartir}
\usepackage{listings}
\usepackage{xcolor}

\lstset{
  basicstyle=\ttfamily\small,
  breaklines=true,
  columns=fullflexible,
  keepspaces=true,
  frame=single,
  showstringspaces=false
}

\begin{document}

\title{Basic Properties of System F}

\author{Runbang Yan}
\email{rbyan25@stu.pku.edu.cn}
\affiliation{%
  \institution{School of Computer Science, Peking University}
  \city{Beijing}
  \country{China}
}

\renewcommand{\shortauthors}{Yan}

\maketitle

\section{Preservation Theorem}

\textbf{Theorem.}
If $\Gamma \vdash t : T$ and $t \longrightarrow t'$, then $\Gamma \vdash t' : T$.

\textbf{Proof.}
We prove the preservation theorem for System F by induction on the typing derivation of $\Gamma \vdash t : T$.

The cases inherited from the simply typed lambda-calculus are standard. Therefore, we only consider the newly added cases for type abstraction and type application.

\begin{itemize}
    \item \textbf{Case \textsc{T-TABS}.}

    Suppose the last rule used in the typing derivation is \textsc{T-TABS}. Then $t = \lambda X.\,t_2$, and the typing derivation has the form
    \[
    \inferrule*
    {
      \Gamma, X \vdash t_2 : T_2
    }
    {
      \Gamma \vdash \lambda X.\,t_2 : \forall X.\,T_2
    }.
    \]

    Since a type abstraction is a value, it cannot take a step. This contradicts the premise that $t \longrightarrow t'$. Therefore, this case is vacuously true.

    \item \textbf{Case \textsc{T-TAPP}.}

    Suppose the last rule used in the typing derivation is \textsc{T-TAPP}. Then $t = t_1\,[T_2]$, and the typing derivation has the form
    \[
    \inferrule*
    {
      \Gamma \vdash t_1 : \forall X.\,T_1
    }
    {
      \Gamma \vdash t_1\,[T_2] : [X \mapsto T_2]T_1
    }.
    \]
    Hence $T = [X \mapsto T_2]T_1$.

    There are two possible ways in which $t_1\,[T_2]$ can take a step.

    \begin{enumerate}
        \item Suppose $t_1 \longrightarrow t_1'$. Then, by the evaluation rule for type application, $t_1\,[T_2] \longrightarrow t_1'\,[T_2]$.

        By the induction hypothesis applied to $\Gamma \vdash t_1 : \forall X.\,T_1$ and $t_1 \longrightarrow t_1'$, we have $\Gamma \vdash t_1' : \forall X.\,T_1$.

        Therefore, by \textsc{T-TAPP}, $\Gamma \vdash t_1'\,[T_2] : [X \mapsto T_2]T_1$. Since $T = [X \mapsto T_2]T_1$, we get $\Gamma \vdash t_1'\,[T_2] : T$.

        \item Suppose $t_1 = \lambda X.\,t_{12}$. Then, by the evaluation rule for type-level beta-reduction,
        \[
        (\lambda X.\,t_{12})\,[T_2]
        \longrightarrow
        [X \mapsto T_2]t_{12}.
        \]

        Since $\Gamma \vdash \lambda X.\,t_{12} : \forall X.\,T_1$, by inversion of \textsc{T-TABS}, we have $\Gamma, X \vdash t_{12} : T_1$.

        By the type substitution lemma, we obtain $\Gamma \vdash [X \mapsto T_2]t_{12} : [X \mapsto T_2]T_1$.

        Since $T = [X \mapsto T_2]T_1$, we conclude that $\Gamma \vdash [X \mapsto T_2]t_{12} : T$.
    \end{enumerate}
\end{itemize}

Thus all new cases preserve types. Together with the standard simply typed lambda-calculus cases, preservation holds for System F.
\[
\Box
\]

\section{Progress Theorem}

\textbf{Theorem.}
If $\emptyset \vdash t : T$, then either $t$ is a value, or there exists some term $t'$ such that $t \longrightarrow t'$.

\textbf{Proof.}
We prove the progress theorem for System F by induction on the typing derivation of $\emptyset \vdash t : T$.

The cases inherited from the simply typed lambda-calculus are standard. Therefore, we only consider the newly added cases for type abstraction and type application.

\begin{itemize}
    \item \textbf{Case \textsc{T-TABS}.}

    Suppose the last rule used in the typing derivation is \textsc{T-TABS}. Then $t = \lambda X.\,t_2$.

    By definition, a type abstraction is a value. Therefore, progress holds in this case.

    \item \textbf{Case \textsc{T-TAPP}.}

    Suppose the last rule used in the typing derivation is \textsc{T-TAPP}. Then $t = t_1\,[T_2]$, and the typing derivation has the form
    \[
    \inferrule*
    {
      \emptyset \vdash t_1 : \forall X.\,T_1
    }
    {
      \emptyset \vdash t_1\,[T_2] : [X \mapsto T_2]T_1
    }.
    \]

    By the induction hypothesis applied to $\emptyset \vdash t_1 : \forall X.\,T_1$, either $t_1$ is a value, or there exists some term $t_1'$ such that $t_1 \longrightarrow t_1'$.

    \begin{enumerate}
        \item Suppose $t_1 \longrightarrow t_1'$. Then, by the evaluation rule for type application, $t_1\,[T_2] \longrightarrow t_1'\,[T_2]$. Therefore, $t$ can take a step.

        \item Suppose $t_1$ is a value.

        Since $\emptyset \vdash t_1 : \forall X.\,T_1$, by the canonical forms lemma for universal types, there exists a term $t_{12}$ such that $t_1 = \lambda X.\,t_{12}$.

        Hence $t = (\lambda X.\,t_{12})\,[T_2]$. By type-level beta-reduction,
        \[
        (\lambda X.\,t_{12})\,[T_2]
        \longrightarrow
        [X \mapsto T_2]t_{12}.
        \]

        Therefore, $t$ can take a step.
    \end{enumerate}
\end{itemize}

Thus all new cases satisfy progress. Together with the standard simply typed lambda-calculus cases, progress holds for System F.
\[
\Box
\]

\end{document}
